\section{Pregunta 3: Aplicación de DDD y Bounded Context en un dominio complejo}

% Responsable: Alfonso Tejeda (@alfonsotejeda)

\subsection*{El síntoma: un término, cuatro modelos}

El escenario planteado describe un sistema en el que el término \emph{póliza} designa cuatro
realidades distintas según quién lo pronuncie: para el actuario es un contrato con cláusulas y
coberturas; para el departamento de siniestros, el expediente contra el que se presenta una
reclamación; para facturación, una fuente de ingresos recurrentes; y para el ajustador que evalúa
los daños en el lugar del accidente, apenas un número de referencia.

Conviene empezar por reconocer de qué clase de problema se trata. No es una imprecisión
terminológica que pueda resolverse acordando una definición común entre departamentos, sino una
propiedad del dominio: cada área necesita de la póliza precisamente aquello que le permite ejecutar
su trabajo, y las reglas de negocio que debe hacer cumplir sobre ella difieren entre sí.
\textcite{evans2003} sostiene que el intento de construir un modelo único que unifique un dominio
extenso conduce a uno de dos resultados, ambos indeseables: un modelo tan genérico que deja de
expresar las reglas del negocio, o un modelo que acumula las preocupaciones de todas las áreas y
que, en consecuencia, nadie puede modificar sin afectar a las demás. La alternativa que propone no
consiste en unificar mejor, sino en renunciar deliberadamente a la unificación total y delimitar de
forma explícita el ámbito de validez de cada modelo.

\subsection*{División en contextos delimitados}

Un \emph{Bounded Context} o contexto delimitado es la frontera explícita dentro de la cual un
modelo de dominio concreto está definido y resulta aplicable, y dentro de la cual su terminología
---el \emph{Lenguaje Ubicuo}--- es consistente y carece de ambigüedad \parencite{evans2003}. La
frontera es simultáneamente lingüística, organizativa y técnica: delimita un vocabulario, el equipo
que lo emplea y el código y el esquema de datos que lo materializan.
\textcite{richardsford2020} destacan además su papel como criterio para decidir qué elementos
pertenecen juntos al establecer la granularidad de los componentes de un sistema.

El propio enunciado proporciona el criterio de división: allí donde un mismo término sostiene
significados incompatibles, pasa una frontera. Se proponen por tanto cuatro contextos delimitados.

\begin{center}
\begin{tabular}{p{3.1cm}p{4.4cm}p{5.4cm}}
\hline
\textbf{Contexto} & \textbf{Qué es la póliza} & \textbf{Reglas que hace cumplir} \\
\hline
Suscripción & Contrato con cláusulas, coberturas y vigencia & Tarificación, evaluación del riesgo, condiciones de emisión y renovación \\
\hline
Siniestros & Expediente contra el que se reclama & Cobertura del hecho declarado, estados de la reclamación, plazos \\
\hline
Facturación de primas & Fuente de ingresos recurrentes & Cuotas, recibos, vencimientos, tratamiento del impago \\
\hline
Peritaje en campo & Número de referencia de una inspección & Valoración de daños, evidencia fotográfica, cierre del informe \\
\hline
\end{tabular}
\end{center}

Cada contexto mantiene su propio modelo de póliza, con los atributos y el comportamiento que sus
reglas requieren y sin ninguno de los que no requieren. Lo que comparten no es una entidad, sino
una \emph{identidad}: el identificador de la póliza, que permite correlacionar la información de los
cuatro modelos sin que ninguno dependa de la estructura interna de los otros.

Cabe precisar que un contexto delimitado no es un módulo, una capa ni un microservicio. Es una
frontera semántica y, aunque puede coincidir con una unidad de despliegue, se trata de decisiones de
distinto nivel: la primera responde a la coherencia del modelo; la segunda, a requisitos de
escalabilidad, autonomía de los equipos y operación. Identificar ambas de forma automática es uno de
los errores que \textcite{richardsford2020} advierten al discutir la granularidad de los servicios.

\subsection*{Protección del modelo de Póliza}

Delimitar los contextos no basta: hace falta un mecanismo que impida que los detalles de un contexto
se filtren hacia el modelo de otro. Ese mecanismo descansa sobre dos piezas que operan de forma
conjunta.

La primera es la ya mencionada ausencia de una entidad compartida. En el contexto de Suscripción,
\texttt{Póliza} es una \emph{raíz de agregado} \parencite{evans2003}: sus entidades y objetos de
valor internos ---cláusulas, coberturas, condiciones particulares--- solo son alcanzables a través
de la raíz, que es la responsable de hacer cumplir las invariantes del contrato. Ningún componente
externo al agregado conserva una referencia directa a una cláusula suelta. Facturación no manipula
esa estructura: mantiene su propia \texttt{CuentaDePrima}, que referencia la póliza únicamente por
su identificador. Siniestros hace lo mismo con su \texttt{Expediente}, y Peritaje con su
\texttt{OrdenDeInspección}.

La segunda es el Lenguaje Ubicuo de cada contexto, que actúa como criterio de admisión: un atributo
solo entra en un modelo si el vocabulario de ese contexto lo reclama. Preguntas como
\enquote{¿cuánto se debe de esta póliza?} o \enquote{¿qué fotografías tomó el ajustador?} no
pertenecen al Lenguaje Ubicuo de Suscripción y, por tanto, no justifican atributos en su modelo.

El valor de estas dos restricciones se aprecia mejor por contraste. Supóngase el modelo unificado que
se obtiene al ignorarlas, con una única clase que reúna cláusulas, coberturas, recibos pendientes,
saldo deudor, siniestros abiertos, fotografías de los daños y la ubicación del ajustador. Las
consecuencias son tres, en orden creciente de gravedad. Primera, un cambio en las reglas de
facturación obliga a modificar y redesplegar el componente que sostiene la suscripción, pese a que
las reglas actuariales no han cambiado. Segunda, las invariantes de las distintas áreas entran en
conflicto dentro de una misma clase: para facturación una póliza con recibos impagados sigue
existiendo y debe seguir generando cuotas, mientras que para suscripción puede haber quedado sin
vigencia, y una sola clase no puede sostener ambas reglas a la vez sin recurrir a condicionales que
dependen de quién la invoca. Tercera, y decisiva, la expresión \enquote{póliza activa} deja de tener
un significado único en el código, que es exactamente la pérdida del Lenguaje Ubicuo que la
delimitación de contextos venía a evitar.

Por la misma razón debe descartarse resolver la integración mediante una base de datos compartida
sobre la que los cuatro contextos leen y escriben. Una tabla común de pólizas reintroduce el modelo
unificado por la puerta del esquema relacional: las cuatro áreas vuelven a quedar acopladas a una
misma estructura, y cualquier columna añadida para satisfacer a una de ellas pasa a formar parte del
modelo de las demás. \textcite{evans2003} contempla el patrón \emph{Shared Kernel} para compartir
de forma deliberada un subconjunto reducido del modelo, pero con un coste de coordinación permanente
entre los equipos implicados; en un escenario de cuatro áreas con vocabularios divergentes, ese
coste no se justifica.

\subsection*{Comunicación desacoplada entre los contextos}

Los contextos permanecen separados pero no aislados: la emisión de una póliza debe producir cuotas,
y la declaración de un siniestro debe originar una inspección. \textcite{evans2003} propone
documentar esas relaciones de forma explícita en un \emph{Context Map}, acompañado de un catálogo de
patrones de integración. Tres de ellos resultan pertinentes aquí.

\begin{itemize}
  \item \textbf{Capa de Anticorrupción} (\emph{Anticorruption Layer}). Es el patrón indicado para
  consumir un modelo ajeno que no se controla, como un sistema de facturación de primas preexistente.
  La capa traduce el modelo externo al propio en la frontera, de modo que el vocabulario y las
  decisiones de diseño del sistema ajeno no penetren en el modelo de dominio.

  \item \textbf{Open Host Service} con \textbf{Published Language}. Suscripción atiende a varios
  consumidores, por lo que en lugar de exponer su modelo interno publica un contrato estable y
  versionado, expresado en un lenguaje de intercambio acordado. Así puede evolucionar su modelo
  interno sin romper a sus consumidores, siempre que respete el contrato publicado.

  \item \textbf{Customer/Supplier}. Describe la relación entre Siniestros y Peritaje en campo: el
  segundo provee al primero y ambos pertenecen a la misma organización, de modo que las necesidades
  del consumidor pueden negociarse y planificarse en lugar de simplemente padecerse.
\end{itemize}

Resta justificar por qué la comunicación debe ser predominantemente asíncrona, y el argumento no es
de oportunidad tecnológica sino una consecuencia del modelo. La regla de consistencia de los
agregados establece que una transacción modifica un solo agregado \parencite{evans2003}; en
consecuencia, en el momento en que una operación de negocio atraviesa la frontera de un contexto, la
consistencia transaccional deja de estar disponible por construcción. Lo que queda es la
publicación de hechos del dominio ---\emph{póliza emitida}, \emph{siniestro declarado},
\emph{prima impagada}--- mediante mensajería, y la reacción autónoma de cada contexto, que mantiene
su propia proyección de aquello que necesita saber del resto.

El compromiso correspondiente debe declararse con honestidad. La integración asíncrona aporta
autonomía operativa y tolerancia a la indisponibilidad temporal de un contexto, pero a cambio
introduce consistencia eventual: existen ventanas durante las cuales Facturación aún no conoce una
póliza ya emitida. Sostener ese diseño exige razonar sobre el orden y la posible duplicación de los
mensajes, y aceptar que ciertas preguntas solo admiten respuesta aproximada hasta que el sistema
converge; se trata de uno de los compromisos característicos de la integración distribuida que
analizan \textcite{richardsford2020}.

\subsection*{Arquitectura orientada al dominio y Adaptador de Servicio}

La solución anterior describe \emph{dónde} están las fronteras, pero no \emph{cómo} se sostienen
dentro de cada contexto. Esa es la aportación de las arquitecturas orientadas al dominio.
La Arquitectura Cebolla organiza el sistema en capas concéntricas con el Modelo de Dominio en el
centro, de manera que todas las dependencias apuntan hacia adentro y las capas externas implementan
interfaces definidas en las internas \parencite{palermo2008}. La Arquitectura Limpia formula el
mismo compromiso como \emph{Regla de Dependencia}: las dependencias del código fuente solo pueden
apuntar hacia adentro, hacia las políticas de mayor nivel \parencite{martin2017}.

La relación entre ambos cuerpos de ideas es de complementariedad y opera en niveles distintos: DDD
determina \emph{qué} ocupa el centro ---el modelo de dominio de ese contexto, con su Lenguaje
Ubicuo--- mientras que la arquitectura determina \emph{cómo} se lo protege, fijando la dirección de
las dependencias. De ahí se sigue una decisión concreta de diseño: no existe una única cebolla para
todo el sistema, sino que cada contexto delimitado recibe su propia instancia de la arquitectura,
con su propio modelo de dominio en el centro y su propia infraestructura en el exterior.

Queda por ubicar la traducción entre contextos. La Guía de Arquitectura N-Capas orientada al Dominio
sitúa en la capa de Infraestructura unos componentes denominados \emph{Adaptadores de Servicio}, cuya
responsabilidad es encapsular la invocación de servicios externos y traducir los datos que estos
devuelven al modelo propio, de modo que ni el protocolo ni el vocabulario ajenos asciendan hacia las
capas de dominio y aplicación \parencite{delatorre2010}.

Esto permite cerrar el razonamiento: el Adaptador de Servicio es el lugar donde se materializa la
Capa de Anticorrupción. El núcleo de dominio declara la interfaz que necesita ---por ejemplo, un
\texttt{ServicioDeFacturación} expresado en los términos de su propio Lenguaje Ubicuo--- y el
adaptador, residente en la capa de infraestructura, la implementa contra el protocolo concreto del
sistema de primas y efectúa la traducción entre ambos modelos. La interfaz pertenece a la capa
interna y la implementación a la externa, con lo que la dependencia del código fuente apunta hacia
adentro y la Regla de Dependencia queda satisfecha. Es la misma construcción que
\textcite{cockburn2005} describe como puertos y adaptadores, con el puerto en poder del interior.

Dos consecuencias prácticas se desprenden de esta disposición. La primera es la testabilidad: al
depender el dominio de una abstracción propia y no del cliente concreto del sistema externo, su
lógica puede ejercitarse con una implementación sustituta de la interfaz, sin invocar el servicio de
primas ni la biblioteca de comunicación que lo consume. La segunda es la capacidad de evolución:
sustituir el sistema de facturación, o su protocolo, afecta a un único adaptador en la
infraestructura y deja intacto el modelo de dominio, que es la porción del sistema cuya estabilidad
interesa preservar.

\subsection*{Síntesis}

La solución propuesta se articula en tres niveles que responden a preguntas distintas y que deben
sostenerse simultáneamente. En el nivel estratégico, DDD determina dónde situar las fronteras:
cuatro contextos delimitados, deducidos de los significados incompatibles que el dominio asigna al
término póliza, cada uno con su Lenguaje Ubicuo y su propio modelo. En el nivel estructural, las
arquitecturas Cebolla y Limpia fijan la dirección de las dependencias dentro de cada contexto y
mantienen el modelo de dominio libre de detalles de infraestructura. En el nivel de implementación,
el Adaptador de Servicio concreta la traducción en la frontera y da forma a la Capa de
Anticorrupción, de manera que el modelo protegido en los dos niveles anteriores sobreviva al
contacto con sistemas que la organización no controla.
